\subsection{\label{sec.det.comb.lgv}The Lindstr\"{o}m--Gessel--Viennot lemma}

We have so far mostly been discussing algebraic properties of determinants, if
often from a combinatorial point of view. We shall now see a situation where
determinants naturally appear in combinatorics.

Namely, we will see an application of determinants to lattice path enumeration
-- i.e., to the counting of paths on a certain infinite digraph called the
\emph{integer lattice}. A survey of this subject can be found in
\cite{Kratte17}; we will restrict ourselves to one of the most accessible
highlights: the \emph{Lindstr\"{o}m--Gessel--Viennot lemma} (for short:
\emph{LGV lemma}). This \textquotedblleft lemma\textquotedblright\ is by now
so classical and popular that it is often deservedly called the
\emph{Lindstr\"{o}m--Gessel--Viennot theorem}. Alternate treatments of this
theorem can be found in \cite[\S 2.5]{Sagan19}, in \cite[\S 2.7]{Stanley-EC1}
and in \cite[\S 2]{GesVie89}. Some applications predating the general
statement of the theorem can be found in \cite{GesVie85}, and a surprising
recent generalization in \cite[Theorem 2.5]{Talask12}.

\subsubsection{Definitions}

We have already seen lattice paths (and even counted them in Subsection
\ref{subsec.qbin.motiv}). We shall now introduce them formally and study them
in greater depth.

\Needspace{20pc}

\begin{convention}
\label{conv.lgv.digraph}\textbf{(a)} Recall that \textquotedblleft%
\emph{digraph}\textquotedblright\ means \textquotedblleft directed
graph\textquotedblright, i.e., a graph whose edges are directed (and are
called \emph{arcs}). Against a widespread convention, we will allow our
digraphs to be infinite (i.e., to have infinitely many vertices and arcs).
\medskip

\textbf{(b)} A digraph $D$ will be called \emph{path-finite} if it has the
property that for any two vertices $u$ and $v$, there are only finitely many
paths from $u$ to $v$. (Thus, in particular, such paths can be counted.)
\medskip

\textbf{(c)} A digraph $D$ will be called \emph{acyclic} if it has no directed
cycles. For example, the digraph $%
%TCIMACRO{\TeXButton{tikz digraph}{\raisebox{-3.3pc}{
%% [tikz figure removed]
%}}}%
%BeginExpansion
\raisebox{-3.3pc}{
% [tikz figure removed]
}%
%EndExpansion
$ is acyclic, whereas the digraph $%
%TCIMACRO{\TeXButton{tikz digraph}{\raisebox{-3.3pc}{
%% [tikz figure removed]
%}}}%
%BeginExpansion
\raisebox{-3.3pc}{
% [tikz figure removed]
}%
%EndExpansion
$ is not. \medskip

\textbf{(d)} A \emph{simple digraph} $D$ means a digraph whose arcs are merely
pairs of distinct vertices (i.e., each arc is a pair $\left(  u,v\right)  $ of
two vertices $u$ and $v$ with $u\neq v$).
\end{convention}

We note that a path may contain $0$ arcs (in which case its starting and
ending point are identical).

\begin{definition}
\label{def.lgv.lattice}We consider the infinite simple digraph with vertex set
$\mathbb{Z}^{2}$ (so the vertices are pairs of integers) and arcs%
\begin{equation}
\left(  i,j\right)  \rightarrow\left(  i+1,j\right)
\ \ \ \ \ \ \ \ \ \ \text{for all }\left(  i,j\right)  \in\mathbb{Z}^{2}
\label{eq.def.lgv.lattice.east}%
\end{equation}
and%
\begin{equation}
\left(  i,j\right)  \rightarrow\left(  i,j+1\right)
\ \ \ \ \ \ \ \ \ \ \text{for all }\left(  i,j\right)  \in\mathbb{Z}^{2}.
\label{eq.def.lgv.lattice.north}%
\end{equation}
The arcs of the form (\ref{eq.def.lgv.lattice.east}) are called
\textquotedblleft\emph{east-steps}\textquotedblright\ or \textquotedblleft%
\emph{right-steps}\textquotedblright; the arcs of the form
(\ref{eq.def.lgv.lattice.north}) are called \textquotedblleft%
\emph{north-steps}\textquotedblright\ or \textquotedblleft\emph{up-steps}%
\textquotedblright.

The vertices of this digraph will be called \emph{lattice points} or
\emph{grid points} or simply \emph{points}. They will be represented as points
in the Cartesian plane (in the usual way: the vertex $\left(  i,j\right)
\in\mathbb{Z}^{2}$ is represented as the point with x-coordinate $i$ and
y-coordinate $j$).

The entire digraph will be denoted by $\mathbb{Z}^{2}$ and called the
\emph{integer lattice} or \emph{integer grid} (or, to be short, just
\emph{lattice} or \emph{grid}). Here is a picture of a small part of this
digraph $\mathbb{Z}^{2}$:%
\[%
% [tikz figure removed]
%BeginExpansion
% [tikz figure removed]%
%EndExpansion
\ \ .
\]
Formally speaking, this path is the $9$-tuple%
\[
\left(  \left(  0,0\right)  ,\left(  0,1\right)  ,\left(  1,1\right)  ,\left(
2,1\right)  ,\left(  3,1\right)  ,\left(  3,2\right)  ,\left(  4,2\right)
,\left(  4,3\right)  ,\left(  5,3\right)  \right)  .
\]
Its step sequence\ (i.e., the sequence of the directions of its steps) is
$URRRURUR$ (meaning that its first step is an up-step, its second step is a
right-step, its third step is a right-step, and so on).
\end{example}

Clearly, any path is uniquely determined by its starting point and its step sequence.

Note that we are considering one of the simplest possible notions of a lattice
path here. In more advanced texts, the word \textquotedblleft lattice
path\textquotedblright\ is often used for paths in digraphs more complicated
than $\mathbb{Z}^{2}$ (for instance, a digraph with the same vertex set
$\mathbb{Z}^{2}$ but allowing steps in all four directions). However, the
digraph we are considering is perhaps the most useful for algebraic combinatorics.

\subsubsection{Counting paths from $\left(  a,b\right)  $ to $\left(
c,d\right)  $}

In Subsection \ref{subsec.qbin.motiv}, we have counted the lattice paths from
$\left(  0,0\right)  $ to $\left(  6,4\right)  $ that begin with an east-step
and end with a north-step. These are, of course, in bijection with the paths
from $\left(  1,0\right)  $ to $\left(  6,3\right)  $ (since the first and
last step are predetermined and thus can be ignored). Let us now generalize
this by counting paths between any two lattice points:

\begin{proposition}
\label{prop.lgv.1-paths.ct}Let $\left(  a,b\right)  \in\mathbb{Z}^{2}$ and
$\left(  c,d\right)  \in\mathbb{Z}^{2}$ be two points. Then,%
\[
\left(  \text{\# of paths from }\left(  a,b\right)  \text{ to }\left(
c,d\right)  \right)  =%
\begin{cases}
\dbinom{c+d-a-b}{c-a}, & \text{if }c+d\geq a+b;\\
0, & \text{if }c+d<a+b.
\end{cases}
\]

\end{proposition}

\begin{proof}
[Proof of Proposition \ref{prop.lgv.1-paths.ct}.]This is just a formalization
(and generalization) of the reasoning we used in Subsection
\ref{subsec.qbin.motiv}.

We shall first show the following two observations:

\begin{statement}
\textit{Observation 1:} Each path from $\left(  a,b\right)  $ to $\left(
c,d\right)  $ has exactly $c+d-a-b$ steps\footnote{A \textquotedblleft
step\textquotedblright\ of a path means an arc of this path.}.
\end{statement}

\begin{statement}
\textit{Observation 2:} Each path from $\left(  a,b\right)  $ to $\left(
c,d\right)  $ has exactly $c-a$ east-steps.
\end{statement}

\begin{fineprint}
[\textit{Proof of Observation 1:} We define the \emph{coordinate sum} of a
point $\left(  x,y\right)  \in\mathbb{Z}^{2}$ to be $x+y$. We shall denote
this coordinate sum by $\operatorname*{cs}\left(  x,y\right)  $. We observe
that the coordinate sum of a point increases by exactly $1$ along each arc of
$\mathbb{Z}^{2}$: That is, if $u\rightarrow v$ is an arc of $\mathbb{Z}^{2}$,
then%
\begin{equation}
\operatorname*{cs}\left(  v\right)  -\operatorname*{cs}\left(  u\right)  =1.
\label{pf.prop.lgv.1-paths.ct.o1.pf.1}%
\end{equation}
(This is because we can write $u$ in the form $u=\left(  i,j\right)  $ and
then must have either $v=\left(  i+1,j\right)  $ or $v=\left(  i,j+1\right)
$; but this entails $\operatorname*{cs}\left(  v\right)  =\underbrace{i+j}%
_{=\operatorname*{cs}\left(  u\right)  }+\,1=\operatorname*{cs}\left(
u\right)  +1$ in either case.)

Let $\left(  v_{0},v_{1},\ldots,v_{n}\right)  $ be a path from $\left(
a,b\right)  $ to $\left(  c,d\right)  $. Thus, $v_{0}=\left(  a,b\right)  $
and $v_{n}=\left(  c,d\right)  $. Moreover, for each $i\in\left[  n\right]  $,
we know that $v_{i-1}\rightarrow v_{i}$ is an arc of $\mathbb{Z}^{2}$, and
thus we have%
\[
\operatorname*{cs}\left(  v_{i}\right)  -\operatorname*{cs}\left(
v_{i-1}\right)  =1
\]
(by (\ref{pf.prop.lgv.1-paths.ct.o1.pf.1}), applied to $u=v_{i-1}$ and
$v=v_{i}$). Summing these equalities over all $i\in\left[  n\right]  $, we
obtain%
\[
\sum_{i=1}^{n}\left(  \operatorname*{cs}\left(  v_{i}\right)
-\operatorname*{cs}\left(  v_{i-1}\right)  \right)  =\sum_{i=1}^{n}1=n.
\]
Hence,%
\begin{align*}
n  &  =\sum_{i=1}^{n}\left(  \operatorname*{cs}\left(  v_{i}\right)
-\operatorname*{cs}\left(  v_{i-1}\right)  \right)  =\operatorname*{cs}%
\underbrace{\left(  v_{n}\right)  }_{=\left(  c,d\right)  }-\operatorname*{cs}%
\underbrace{\left(  v_{0}\right)  }_{=\left(  a,b\right)  }%
\ \ \ \ \ \ \ \ \ \ \left(  \text{by the telescope principle}\right) \\
&  =\underbrace{\operatorname*{cs}\left(  c,d\right)  }_{=c+d}%
-\underbrace{\operatorname*{cs}\left(  a,b\right)  }_{=a+b}=c+d-\left(
a+b\right)  =c+d-a-b.
\end{align*}
In other words, the path $\left(  v_{0},v_{1},\ldots,v_{n}\right)  $ has
exactly $c+d-a-b$ steps (since this path clearly has $n$ steps).

Forget that we fixed $\left(  v_{0},v_{1},\ldots,v_{n}\right)  $. We thus have
shown that each path $\left(  v_{0},v_{1},\ldots,v_{n}\right)  $ from $\left(
a,b\right)  $ to $\left(  c,d\right)  $ has exactly $c+d-a-b$ steps. This
proves Observation 1.] \medskip
\end{fineprint}

\begin{fineprint}
[\textit{Proof of Observation 2:} For any point $v\in\mathbb{Z}^{2}$, we
define $\operatorname*{x}\left(  v\right)  $ to be the x-coordinate of $v$.
(Thus, $\operatorname*{x}\left(  x,y\right)  =x$ for each $\left(  x,y\right)
\in\mathbb{Z}^{2}$.)

Obviously, the x-coordinate of a point increases by exactly $1$ along each
east-step and stays unchanged along each north-step: That is, if $u\rightarrow
v$ is an arc of $\mathbb{Z}^{2}$, then%
\begin{equation}
\operatorname*{x}\left(  v\right)  -\operatorname*{x}\left(  u\right)  =%
\begin{cases}
1, & \text{if }u\rightarrow v\text{ is an east-step;}\\
0, & \text{if }u\rightarrow v\text{ is a north-step.}%
\end{cases}
\label{pf.prop.lgv.1-paths.ct.o2.pf.1}%
\end{equation}

\end{fineprint}

Let $\left(  v_{0},v_{1},\ldots,v_{n}\right)  $ be a path from $\left(
a,b\right)  $ to $\left(  c,d\right)  $. Thus, $v_{0}=\left(  a,b\right)  $
and $v_{n}=\left(  c,d\right)  $. Moreover, for each $i\in\left[  n\right]  $,
we know that $v_{i-1}\rightarrow v_{i}$ is an arc of $\mathbb{Z}^{2}$, and
thus we have%
\[
\operatorname*{x}\left(  v_{i}\right)  -\operatorname*{x}\left(
v_{i-1}\right)  =%
\begin{cases}
1, & \text{if }v_{i-1}\rightarrow v_{i}\text{ is an east-step;}\\
0, & \text{if }v_{i-1}\rightarrow v_{i}\text{ is a north-step}%
\end{cases}
\]
(by (\ref{pf.prop.lgv.1-paths.ct.o2.pf.1}), applied to $u=v_{i-1}$ and
$v=v_{i}$). Summing these equalities over all $i\in\left[  n\right]  $, we
obtain%
\begin{align*}
\sum_{i=1}^{n}\left(  \operatorname*{x}\left(  v_{i}\right)
-\operatorname*{x}\left(  v_{i-1}\right)  \right)   &  =\sum_{i=1}^{n}%
\begin{cases}
1, & \text{if }v_{i-1}\rightarrow v_{i}\text{ is an east-step;}\\
0, & \text{if }v_{i-1}\rightarrow v_{i}\text{ is a north-step}%
\end{cases}
\\
&  =\left(  \text{\# of }i\in\left[  n\right]  \text{ such that }%
v_{i-1}\rightarrow v_{i}\text{ is an east-step}\right) \\
&  =\left(  \text{\# of east-steps in the path }\left(  v_{0},v_{1}%
,\ldots,v_{n}\right)  \right)  .
\end{align*}
Hence,%
\begin{align*}
&  \left(  \text{\# of east-steps in the path }\left(  v_{0},v_{1}%
,\ldots,v_{n}\right)  \right) \\
&  =\sum_{i=1}^{n}\left(  \operatorname*{x}\left(  v_{i}\right)
-\operatorname*{x}\left(  v_{i-1}\right)  \right)  =\operatorname*{x}%
\underbrace{\left(  v_{n}\right)  }_{=\left(  c,d\right)  }-\operatorname*{x}%
\underbrace{\left(  v_{0}\right)  }_{=\left(  a,b\right)  }%
\ \ \ \ \ \ \ \ \ \ \left(  \text{by the telescope principle}\right) \\
&  =\underbrace{\operatorname*{x}\left(  c,d\right)  }_{=c}%
-\underbrace{\operatorname*{x}\left(  a,b\right)  }_{=a}=c-a.
\end{align*}
In other words, the path $\left(  v_{0},v_{1},\ldots,v_{n}\right)  $ has
exactly $c-a$ east-steps.

Forget that we fixed $\left(  v_{0},v_{1},\ldots,v_{n}\right)  $. We thus have
shown that each path $\left(  v_{0},v_{1},\ldots,v_{n}\right)  $ from $\left(
a,b\right)  $ to $\left(  c,d\right)  $ has exactly $c-a$ east-steps. This
proves Observation 2.] \medskip

Observation 1 immediately shows that no path from $\left(  a,b\right)  $ to
$\left(  c,d\right)  $ exists when $c+d-a-b<0$. In other words, no path from
$\left(  a,b\right)  $ to $\left(  c,d\right)  $ exists when $c+d<a+b$. In
other words, $\left(  \text{\# of paths from }\left(  a,b\right)  \text{ to
}\left(  c,d\right)  \right)  =0$ when $c+d<a+b$. This proves Proposition
\ref{prop.lgv.1-paths.ct} in the case when $c+d<a+b$. Hence, for the rest of
this proof of Proposition \ref{prop.lgv.1-paths.ct}, we WLOG assume that
$c+d\geq a+b$. Thus, $c+d-a-b\geq0$, so that $c+d-a-b\in\mathbb{N}$.

Observations 1 and 2 have a sort of (common) converse:

\begin{statement}
\textit{Observation 3:} Let $p$ be a path that starts at the point $\left(
a,b\right)  $ and has exactly $c+d-a-b$ steps. Assume that exactly $c-a$ of
these steps are east-steps. Then, this path $p$ ends at $\left(  c,d\right)  $.
\end{statement}

\begin{fineprint}
[\textit{Proof of Observation 3:} Let $\left(  c^{\prime},d^{\prime}\right)  $
be the point at which this path $p$ ends. Then, we can apply Observation 1 to
$\left(  c^{\prime},d^{\prime}\right)  $ instead of $\left(  c,d\right)  $,
and hence conclude that this path has exactly $c^{\prime}+d^{\prime}-a-b$
steps. Since we already know that this path has exactly $c+d-a-b$ steps, we
therefore conclude that $c^{\prime}+d^{\prime}-a-b=c+d-a-b$. In other words,
$c^{\prime}+d^{\prime}=c+d$. Similarly, using Observation 2, we can find that
$c^{\prime}=c$. Subtracting this equality from $c^{\prime}+d^{\prime}=c+d$, we
obtain $d^{\prime}=d$. Combining $c^{\prime}=c$ with $d^{\prime}=d$, we find
$\left(  c^{\prime},d^{\prime}\right)  =\left(  c,d\right)  $. Thus, our path
$p$ ends at $\left(  c,d\right)  $ (since we know that it ends at $\left(
c^{\prime},d^{\prime}\right)  $). This proves Observation 3.] \medskip
\end{fineprint}

Now, combining Observations 1, 2 and 3, we see that the paths from $\left(
a,b\right)  $ to $\left(  c,d\right)  $ are precisely the paths that start at
$\left(  a,b\right)  $ and have exactly $c+d-a-b$ steps and exactly $c-a$
east-steps (among these $c+d-a-b$ steps). Such a path is therefore uniquely
determined if we know \textbf{which} $c-a$ of its $c+d-a-b$ steps are
east-steps. Thus, specifying such a path is equivalent to specifying a
$\left(  c-a\right)  $-element subset of the $\left(  c+d-a-b\right)
$-element set\footnote{Here we are tacitly using $c+d-a-b\geq0$.} $\left[
c+d-a-b\right]  $. The bijection principle thus yields\footnote{To make this
more formal: We are saying that the map
\begin{align*}
\left\{  \text{paths from }\left(  a,b\right)  \text{ to }\left(  c,d\right)
\right\}   &  \rightarrow\left\{  \left(  c-a\right)  \text{-element subsets
of }\left[  c+d-a-b\right]  \right\}  ,\\
\left(  v_{0},v_{1},\ldots,v_{c+d-a-b}\right)   &  \mapsto\left\{  i\in\left[
c+d-a-b\right]  \ \mid\ \text{the arc }v_{i-1}\rightarrow v_{i}\text{ is an
east-step}\right\}
\end{align*}
is a bijection, and we are applying the bijection principle to this
bijection.}%
\begin{align*}
&  \left(  \text{\# of paths from }\left(  a,b\right)  \text{ to }\left(
c,d\right)  \right) \\
&  =\left(  \text{\# of }\left(  c-a\right)  \text{-element subsets of
}\left[  c+d-a-b\right]  \right) \\
&  =\dbinom{c+d-a-b}{c-a}\ \ \ \ \ \ \ \ \ \ \left(  \text{since }%
c+d-a-b\in\mathbb{N}\right)  .
\end{align*}
This proves Proposition \ref{prop.lgv.1-paths.ct} (since we have assumed that
$c+d\geq a+b$).
\end{proof}

\subsubsection{Path tuples, nipats and ipats}

Now, let us try to count something more interesting: tuples of
non-intersecting paths.

\begin{definition}
\label{def.lgv.path-tups}Let $k\in\mathbb{N}$. \medskip

\textbf{(a)} A $k$\emph{-vertex} means a $k$-tuple of lattice points (i.e., a
$k$-tuple of vertices). For example, $\left(  \left(  1,2\right)  ,\ \left(
4,5\right)  ,\ \left(  7,4\right)  \right)  $ is a $3$-vertex. \medskip

\textbf{(b)} If $\mathbf{A}=\left(  A_{1},A_{2},\ldots,A_{k}\right)  $ is a
$k$-vertex, and if $\sigma\in S_{k}$ is a permutation, then $\sigma\left(
\mathbf{A}\right)  $ shall denote the $k$-vertex $\left(  A_{\sigma\left(
1\right)  },A_{\sigma\left(  2\right)  },\ldots,A_{\sigma\left(  k\right)
}\right)  $. For instance, for the simple transposition $s_{1}\in S_{3}$, we
have $s_{1}\left(  A,B,C\right)  =\left(  B,A,C\right)  $ for any $3$-vertex
$\left(  A,B,C\right)  $. \medskip

\textbf{(c)} If $\mathbf{A}=\left(  A_{1},A_{2},\ldots,A_{k}\right)  $ and
$\mathbf{B}=\left(  B_{1},B_{2},\ldots,B_{k}\right)  $ are two $k$-vertices,
then a \emph{path tuple} from $\mathbf{A}$ to $\mathbf{B}$ means a $k$-tuple
$\left(  p_{1},p_{2},\ldots,p_{k}\right)  $, where each $p_{i}$ is a path from
$A_{i}$ to $B_{i}$. \medskip

\textbf{(d)} A path tuple $\left(  p_{1},p_{2},\ldots,p_{k}\right)  $ is said
to be \emph{non-intersecting} if no two of the paths $p_{1},p_{2},\ldots
,p_{k}$ have any vertex in common. (Visually speaking, this not only forbids
them from crossing each other, but also forbids them from touching or bouncing
off each other, or starting or ending at the same point.)

We shall abbreviate \textquotedblleft non-intersecting path
tuple\textquotedblright\ as \textquotedblleft\emph{nipat}\textquotedblright.
(Historically, the more common abbreviation is \textquotedblleft%
\emph{NILP}\textquotedblright, for \textquotedblleft non-intersecting lattice
paths\textquotedblright, but I prefer \textquotedblleft
nipat\textquotedblright\ as it stresses the tupleness.) \medskip

\textbf{(e)} A path tuple $\left(  p_{1},p_{2},\ldots,p_{k}\right)  $ is said
to be \emph{intersecting} if it is not non-intersecting (i.e., if two of its
paths have a vertex in common).

We shall abbreviate \textquotedblleft intersecting path
tuple\textquotedblright\ as \textquotedblleft\emph{ipat}\textquotedblright.
\end{definition}

\begin{example}
Here are some path tuples for $k=3$: \medskip

\textbf{(a)} The following path tuple is a nipat:%
\[%
%TCIMACRO{\TeXButton{nipat}{% [tikz figure removed]
%}}%
%BeginExpansion
% [tikz figure removed]
%EndExpansion
\ \ .
\]


\textbf{(b)} The following path tuple is an ipat:%
\[%
%TCIMACRO{\TeXButton{ipat}{% [tikz figure removed]}}%
%BeginExpansion
% [tikz figure removed]%
%EndExpansion
\ \ .
\]


\textbf{(c)} The following path tuple is an ipat, too (for several reasons):%
\[%
%TCIMACRO{\TeXButton{ipat}{% [tikz figure removed]}}%
%BeginExpansion
% [tikz figure removed]%
%EndExpansion
\ \ .
\]
(In this tuple, the paths $p_{1}$ and $p_{2}$ even have an arc in common.
Don't let the picture confuse you: The two curved arcs are actually one and
the same arc of $\mathbb{Z}^{2}$ appearing in two paths, not two different arcs.)
\end{example}

\subsubsection{The LGV lemma for two paths}

As I mentioned, we want to count nipats. Here is a first result on the case
$k=2$:

\begin{proposition}
[LGV lemma for two paths]\label{prop.lgv.2paths.count}Let $\left(
A,A^{\prime}\right)  $ and $\left(  B,B^{\prime}\right)  $ be two $2$-vertices
(i.e., let $A,A^{\prime},B,B^{\prime}$ be four lattice points). Then,%
\begin{align*}
&  \det\left(
\begin{array}
[c]{cc}%
\left(  \text{\# of paths from }A\text{ to }B\right)  & \left(  \text{\# of
paths from }A\text{ to }B^{\prime}\right) \\
\left(  \text{\# of paths from }A^{\prime}\text{ to }B\right)  & \left(
\text{\# of paths from }A^{\prime}\text{ to }B^{\prime}\right)
\end{array}
\right) \\
&  =\left(  \text{\# of nipats from }\left(  A,A^{\prime}\right)  \text{ to
}\left(  B,B^{\prime}\right)  \right) \\
&  \ \ \ \ \ \ \ \ \ \ -\left(  \text{\# of nipats from }\left(  A,A^{\prime
}\right)  \text{ to }\left(  B^{\prime},B\right)  \right)  .
\end{align*}

\end{proposition}

\begin{example}
Let $A=\left(  0,0\right)  $ and $A^{\prime}=\left(  1,1\right)  $ and
$B=\left(  2,2\right)  $ and $B^{\prime}=\left(  3,3\right)  $. Then,
Proposition \ref{prop.lgv.2paths.count} says that%
\begin{align*}
\det\left(
\begin{array}
[c]{cc}%
\dbinom{4}{2} & \dbinom{6}{3}\\
\dbinom{2}{1} & \dbinom{4}{2}%
\end{array}
\right)   &  =\left(  \text{\# of nipats from }\left(  A,A^{\prime}\right)
\text{ to }\left(  B,B^{\prime}\right)  \right) \\
&  \ \ \ \ \ \ \ \ \ \ -\left(  \text{\# of nipats from }\left(  A,A^{\prime
}\right)  \text{ to }\left(  B^{\prime},B\right)  \right)  .
\end{align*}
(The matrix entries on the left hand side have been computed using Proposition
\ref{prop.lgv.1-paths.ct}.)

And indeed, this equality is easily verified. There are $2$ nipats from
$\left(  A,A^{\prime}\right)  $ to $\left(  B,B^{\prime}\right)  $, one of
which is
\[%
%TCIMACRO{\TeXButton{nipat}{% [tikz figure removed]}}%
%BeginExpansion
% [tikz figure removed]%
%EndExpansion
\]
while the other is its reflection in the $x=y$ diagonal. There are $6$ nipats
from $\left(  A,A^{\prime}\right)  $ to $\left(  B^{\prime},B\right)  $, three
of which are
\[%
%TCIMACRO{\TeXButton{nipat}{% [tikz figure removed]}}%
%BeginExpansion
% [tikz figure removed]%
%EndExpansion
\ \ \ \ \ \ \ \ \ \
%TCIMACRO{\TeXButton{nipat}{% [tikz figure removed]}}%
%BeginExpansion
% [tikz figure removed]%
%EndExpansion
\ \ \ \ \ \ \ \ \ \
%TCIMACRO{\TeXButton{nipat}{% [tikz figure removed]}}%
%BeginExpansion
% [tikz figure removed]%
%EndExpansion
\]
while the other three are their reflections in the $x=y$ diagonal. The right
hand side of the above equality is thus $2-6=-4$, which is also the left hand side.
\end{example}

\begin{example}
\label{exa.lgv.2paths.2}Let $A=\left(  0,0\right)  $ and $A^{\prime}=\left(
-1,1\right)  $ and $B=\left(  2,2\right)  $ and $B^{\prime}=\left(
0,3\right)  $. Then, the claim of Proposition \ref{prop.lgv.2paths.count}
simplifies, since $\left(  \text{\# of nipats from }\left(  A,A^{\prime
}\right)  \text{ to }\left(  B^{\prime},B\right)  \right)  =0$ in this case.
Here is a picture of the four points that should make this visually clear:
\[%
%TCIMACRO{\TeXButton{just four points}{% [tikz figure removed]}}%
%BeginExpansion
% [tikz figure removed]%
%EndExpansion
\ \ .
\]

\end{example}

\begin{proof}
[Proof of Proposition \ref{prop.lgv.2paths.count}.]We have%
\begin{align}
&  \det\left(
\begin{array}
[c]{cc}%
\left(  \text{\# of paths from }A\text{ to }B\right)  & \left(  \text{\# of
paths from }A\text{ to }B^{\prime}\right) \\
\left(  \text{\# of paths from }A^{\prime}\text{ to }B\right)  & \left(
\text{\# of paths from }A^{\prime}\text{ to }B^{\prime}\right)
\end{array}
\right) \nonumber\\
&  =\underbrace{\left(  \text{\# of paths from }A\text{ to }B\right)
\cdot\left(  \text{\# of paths from }A^{\prime}\text{ to }B^{\prime}\right)
}_{\substack{=\left(  \text{\# of path tuples from }\left(  A,A^{\prime
}\right)  \text{ to }\left(  B,B^{\prime}\right)  \right)  \\\text{(by the
product rule, since a path tuple from }\left(  A,A^{\prime}\right)  \text{ to
}\left(  B,B^{\prime}\right)  \text{ is just}\\\text{a pair consisting of a
path from }A\text{ to }B\text{ and a path from }A^{\prime}\text{ to }%
B^{\prime}\text{)}}}\nonumber\\
&  \ \ \ \ \ \ \ \ \ \ -\underbrace{\left(  \text{\# of paths from }A\text{ to
}B^{\prime}\right)  \cdot\left(  \text{\# of paths from }A^{\prime}\text{ to
}B\right)  }_{\substack{=\left(  \text{\# of path tuples from }\left(
A,A^{\prime}\right)  \text{ to }\left(  B^{\prime},B\right)  \right)
\\\text{(by the product rule, since a path tuple from }\left(  A,A^{\prime
}\right)  \text{ to }\left(  B^{\prime},B\right)  \text{ is just}\\\text{a
pair consisting of a path from }A\text{ to }B^{\prime}\text{ and a path from
}A^{\prime}\text{ to }B\text{)}}}\nonumber\\
&  =\left(  \text{\# of path tuples from }\left(  A,A^{\prime}\right)  \text{
to }\left(  B,B^{\prime}\right)  \right) \nonumber\\
&  \ \ \ \ \ \ \ \ \ \ -\left(  \text{\# of path tuples from }\left(
A,A^{\prime}\right)  \text{ to }\left(  B^{\prime},B\right)  \right)  .
\label{pf.prop.lgv.2paths.count.1}%
\end{align}
We need to show that on the right hand side, all the intersecting path tuples
cancel each other out (so that only the nipats remain).

Our $k$-vertices are $2$-vertices; thus, our path tuples are pairs. Hence,
such a path tuple $\left(  p,p^{\prime}\right)  $ is intersecting if and only
if $p$ and $p^{\prime}$ have a vertex in common. We shall use these common
vertices to define a sign-reversing involution on the intersecting path
tuples. Specifically, we do the following:

Define the set%
\begin{align*}
\mathcal{A}  &  :=\left\{  \text{path tuples from }\left(  A,A^{\prime
}\right)  \text{ to }\left(  B,B^{\prime}\right)  \right\} \\
&  \ \ \ \ \ \ \ \ \ \ \sqcup\left\{  \text{path tuples from }\left(
A,A^{\prime}\right)  \text{ to }\left(  B^{\prime},B\right)  \right\}  .
\end{align*}
Here, the symbol \textquotedblleft$\sqcup$\textquotedblright\ means
\textquotedblleft disjoint union (of sets)\textquotedblright, which is a way
of uniting two sets without removing duplicates (i.e., even if the sets are
not disjoint, we treat them as disjoint for the purpose of the union, and
therefore include two copies of each common element). As a consequence of us
taking the disjoint union, each path tuple in $\mathcal{A}$ \textquotedblleft
remembers\textquotedblright\ whether it comes from the set \newline$\left\{
\text{path tuples from }\left(  A,A^{\prime}\right)  \text{ to }\left(
B,B^{\prime}\right)  \right\}  $ or from the set \newline$\left\{  \text{path
tuples from }\left(  A,A^{\prime}\right)  \text{ to }\left(  B^{\prime
},B\right)  \right\}  $ (and if these two sets have a path tuple in common,
then $\mathcal{A}$ has two copies of it, each of which remembers from which
set it comes). However, in practice, this is barely relevant: Indeed, the only
case in which the sets \newline$\left\{  \text{path tuples from }\left(
A,A^{\prime}\right)  \text{ to }\left(  B,B^{\prime}\right)  \right\}  $ and
$\left\{  \text{path tuples from }\left(  A,A^{\prime}\right)  \text{ to
}\left(  B^{\prime},B\right)  \right\}  $ can fail to be disjoint is the case
when $B=B^{\prime}$; however, in this case, the claim we are proving is
trivial anyway, since there are no nipats\footnote{Indeed, if $p$ and
$p^{\prime}$ are two paths with the same destination, then $p$ and $p^{\prime
}$ automatically have a vertex in common.}, and our matrix has determinant $0$
(since it has two equal columns).

Define a subset $\mathcal{X}$ of $\mathcal{A}$ by%
\[
\mathcal{X}:=\left\{  \text{ipats in }\mathcal{A}\right\}  =\left\{  \left(
p,p^{\prime}\right)  \in\mathcal{A}\ \mid\ p\text{ and }p^{\prime}\text{ have
a vertex in common}\right\}  .
\]
Hence, $\mathcal{A}\setminus\mathcal{X}=\left\{  \text{nipats in }%
\mathcal{A}\right\}  $.

For each $\left(  p,p^{\prime}\right)  \in\mathcal{A}$, we set%
\[
\operatorname*{sign}\left(  p,p^{\prime}\right)  :=%
\begin{cases}
1, & \text{if }\left(  p,p^{\prime}\right)  \text{ is a path tuple from
}\left(  A,A^{\prime}\right)  \text{ to }\left(  B,B^{\prime}\right)  ;\\
-1, & \text{if }\left(  p,p^{\prime}\right)  \text{ is a path tuple from
}\left(  A,A^{\prime}\right)  \text{ to }\left(  B^{\prime},B\right)  .
\end{cases}
\]
(This is well-defined, because each $\left(  p,p^{\prime}\right)
\in\mathcal{A}$ is either a path tuple from $\left(  A,A^{\prime}\right)  $ to
$\left(  B,B^{\prime}\right)  $ or a path tuple from $\left(  A,A^{\prime
}\right)  $ to $\left(  B^{\prime},B\right)  $ but never both at the same
time\footnote{because we took the \textbf{disjoint} union of $\left\{
\text{path tuples from }\left(  A,A^{\prime}\right)  \text{ to }\left(
B,B^{\prime}\right)  \right\}  $ and $\left\{  \text{path tuples from }\left(
A,A^{\prime}\right)  \text{ to }\left(  B^{\prime},B\right)  \right\}  $}.)
Thus,
\begin{align*}
&  \left(  \text{\# of path tuples from }\left(  A,A^{\prime}\right)  \text{
to }\left(  B,B^{\prime}\right)  \right) \\
&  \ \ \ \ \ \ \ \ \ \ -\left(  \text{\# of path tuples from }\left(
A,A^{\prime}\right)  \text{ to }\left(  B^{\prime},B\right)  \right) \\
&  =\sum_{\left(  p,p^{\prime}\right)  \in\mathcal{A}}\operatorname*{sign}%
\left(  p,p^{\prime}\right)
\end{align*}
and%
\begin{align*}
&  \left(  \text{\# of nipats from }\left(  A,A^{\prime}\right)  \text{ to
}\left(  B,B^{\prime}\right)  \right) \\
&  \ \ \ \ \ \ \ \ \ \ -\left(  \text{\# of nipats from }\left(  A,A^{\prime
}\right)  \text{ to }\left(  B^{\prime},B\right)  \right) \\
&  =\sum_{\left(  p,p^{\prime}\right)  \in\mathcal{A}\setminus\mathcal{X}%
}\operatorname*{sign}\left(  p,p^{\prime}\right)  \ \ \ \ \ \ \ \ \ \ \left(
\text{since }\mathcal{A}\setminus\mathcal{X}=\left\{  \text{nipats in
}\mathcal{A}\right\}  \right)  .
\end{align*}
We want to prove that the left hand sides of these two equalities are equal.
Thus, it clearly suffices to show that the right hand sides are equal. By
Lemma \ref{lem.sign.cancel2}, it suffices to find a sign-reversing involution
$f:\mathcal{X}\rightarrow\mathcal{X}$ that has no fixed points.

So let us define our sign-reversing involution $f:\mathcal{X}\rightarrow
\mathcal{X}$. For each path tuple $\left(  p,p^{\prime}\right)  \in
\mathcal{X}$, we define $f\left(  p,p^{\prime}\right)  $ as follows:

\begin{itemize}
\item Since $\left(  p,p^{\prime}\right)  \in\mathcal{X}$, the paths $p$ and
$p^{\prime}$ have a vertex in common. There might be several; let $v$ be the
first one. (The first one on $p$ or the first one on $p^{\prime}$ ? Doesn't
matter, because these are the same thing. Indeed, if the first vertex on $p$
that is contained in $p^{\prime}$ was different from the first vertex on
$p^{\prime}$ that is contained in $p$, then we could obtain a nontrivial
circuit of our digraph by walking from the former vertex to the latter vertex
along $p$ and then back along $p^{\prime}$; but this is impossible, since our
digraph is acyclic.)

We call this vertex $v$ the \emph{first intersection} of $\left(  p,p^{\prime
}\right)  $.

\item Call the part of $p$ that comes after $v$ the \emph{tail} of $p$. Call
the part of $p$ that comes before $v$ the \emph{head}\textbf{ }of $p$.

Call the part of $p^{\prime}$ that comes after $v$ the \emph{tail} of
$p^{\prime}$. Call the part of $p^{\prime}$ that comes before $v$ the
\emph{head}\textbf{ }of $p^{\prime}$.

\item Now, we exchange the tails of the paths $p$ and $p^{\prime}$. That is,
we set%
\begin{align*}
q  &  :=\left(  \text{head of }p\right)  \cup\left(  \text{tail of }p^{\prime
}\right)  \ \ \ \ \ \ \ \ \ \ \text{and}\\
q^{\prime}  &  :=\left(  \text{head of }p^{\prime}\right)  \cup\left(
\text{tail of }p\right)
\end{align*}
(where the symbol \textquotedblleft$\cup$\textquotedblright\ means combining a
path ending at $v$ with a path starting at $v$ in the obvious way), and set
$f\left(  p,p^{\prime}\right)  :=\left(  q,q^{\prime}\right)  $.
\end{itemize}

Thus, we have defined a map $f:\mathcal{X}\rightarrow\mathcal{X}$ (in a
moment, we will explain why it is well-defined). Here is an example:
\[
\text{If }\left(  p,p^{\prime}\right)  =%
%TCIMACRO{\TeXButton{ppprime}{\raisebox{-6pc}{
%% [tikz figure removed]
%}}}%
%BeginExpansion
\raisebox{-6pc}{
% [tikz figure removed]
}%
%EndExpansion
\text{, then }\left(  q,q^{\prime}\right)  =%
%TCIMACRO{\TeXButton{qqprime}{\raisebox{-6pc}{% [tikz figure removed]
%}}}%
%BeginExpansion
\raisebox{-6pc}{% [tikz figure removed]
}%
%EndExpansion
\text{.}%
\]
Here is the same configuration, with the point $v$ marked and with the tails
of the two paths drawn extra-thick:%
\[%
%TCIMACRO{\TeXButton{v and tails}{\raisebox{-6pc}{
%% [tikz figure removed]
%}}}%
%BeginExpansion
\raisebox{-6pc}{
% [tikz figure removed]
}%
%EndExpansion
\ \ .
\]


We note that if $\left(  p,p^{\prime}\right)  \in\mathcal{X}$ is an ipat from
$\left(  A,A^{\prime}\right)  $ to $\left(  B,B^{\prime}\right)  $, then
$f\left(  p,p^{\prime}\right)  =\left(  q,q^{\prime}\right)  $ is an ipat from
$\left(  A,A^{\prime}\right)  $ to $\left(  B^{\prime},B\right)  $ (because by
exchanging the tails of $p$ and $p^{\prime}$, we have caused the two paths to
exchange their destinations as well), and vice versa. Thus, $f\left(
p,p^{\prime}\right)  \in\mathcal{X}$ whenever $\left(  p,p^{\prime}\right)
\in\mathcal{X}$. This shows that the map $f:\mathcal{X}\rightarrow\mathcal{X}$
is well-defined.

Furthermore, let $\left(  p,p^{\prime}\right)  \in\mathcal{X}$ be any ipat,
and let $\left(  q,q^{\prime}\right)  =f\left(  p,p^{\prime}\right)  $. Then,
the vertex $v$ chosen in the definition of $f\left(  p,p^{\prime}\right)  $
(that is, the first intersection of $\left(  p,p^{\prime}\right)  $) is still
the first intersection of $\left(  q,q^{\prime}\right)  $ (because when we
exchange the tails of $p$ and $p^{\prime}$, we do not change their heads, and
thus the resulting paths $q$ and $q^{\prime}$ still do not intersect until
$v$), and therefore this vertex gets chosen again if we apply our map $f$ to
$\left(  q,q^{\prime}\right)  $. As a consequence, $f\left(  q,q^{\prime
}\right)  $ is again $\left(  p,p^{\prime}\right)  $ (since exchanging the
tails of $q$ and $q^{\prime}$ simply undoes the changes incurred when we
exchanged the tails of $p$ and $p^{\prime}$). Thus,
\[
\left(  f\circ f\right)  \left(  p,p^{\prime}\right)  =f\left(
\underbrace{f\left(  p,p^{\prime}\right)  }_{=\left(  q,q^{\prime}\right)
}\right)  =f\left(  q,q^{\prime}\right)  =\left(  p,p^{\prime}\right)  .
\]


Forget that we fixed $\left(  p,p^{\prime}\right)  $. We thus have shown that
$\left(  f\circ f\right)  \left(  p,p^{\prime}\right)  =\left(  p,p^{\prime
}\right)  $ for each $\left(  p,p^{\prime}\right)  \in\mathcal{X}$. Hence,
$f\circ f=\operatorname*{id}$. In other words, $f$ is an involution on
$\mathcal{X}$. Moreover, this involution $f$ is sign-reversing (i.e.,
satisfies $\operatorname*{sign}\left(  f\left(  p,p^{\prime}\right)  \right)
=-\operatorname*{sign}\left(  p,p^{\prime}\right)  $ for any $\left(
p,p^{\prime}\right)  \in\mathcal{X}$)\ \ \ \ \footnote{This is because we have
observed above that if $\left(  p,p^{\prime}\right)  \in\mathcal{X}$ is an
ipat from $\left(  A,A^{\prime}\right)  $ to $\left(  B,B^{\prime}\right)  $,
then $f\left(  p,p^{\prime}\right)  =\left(  q,q^{\prime}\right)  $ is an ipat
from $\left(  A,A^{\prime}\right)  $ to $\left(  B^{\prime},B\right)  $, and
vice versa.}. As a consequence of the latter fact, we see that $f$ has no
fixed points (i.e., that we have $f\left(  p,p^{\prime}\right)  \neq\left(
p,p^{\prime}\right)  $ for any $\left(  p,p^{\prime}\right)  \in\mathcal{X}$).
Hence, Lemma \ref{lem.sign.cancel2} yields%
\begin{equation}
\sum_{\left(  p,p^{\prime}\right)  \in\mathcal{A}}\operatorname*{sign}\left(
p,p^{\prime}\right)  =\sum_{\left(  p,p^{\prime}\right)  \in\mathcal{A}%
\setminus\mathcal{X}}\operatorname*{sign}\left(  p,p^{\prime}\right)  .
\label{pf.prop.lgv.2paths.count.at}%
\end{equation}
As we have explained above, the left hand side of this equality is%
\begin{align*}
&  \left(  \text{\# of path tuples from }\left(  A,A^{\prime}\right)  \text{
to }\left(  B,B^{\prime}\right)  \right) \\
&  \ \ \ \ \ \ \ \ \ \ -\left(  \text{\# of path tuples from }\left(
A,A^{\prime}\right)  \text{ to }\left(  B^{\prime},B\right)  \right)  ,
\end{align*}
whereas its right hand side is%
\begin{align*}
&  \left(  \text{\# of nipats from }\left(  A,A^{\prime}\right)  \text{ to
}\left(  B,B^{\prime}\right)  \right) \\
&  \ \ \ \ \ \ \ \ \ \ -\left(  \text{\# of nipats from }\left(  A,A^{\prime
}\right)  \text{ to }\left(  B^{\prime},B\right)  \right)
\end{align*}
(since $\mathcal{A}\setminus\mathcal{X}=\left\{  \text{nipats in }%
\mathcal{A}\right\}  $). Hence, (\ref{pf.prop.lgv.2paths.count.at}) rewrites
as%
\begin{align*}
&  \left(  \text{\# of path tuples from }\left(  A,A^{\prime}\right)  \text{
to }\left(  B,B^{\prime}\right)  \right) \\
&  \ \ \ \ \ \ \ \ \ \ -\left(  \text{\# of path tuples from }\left(
A,A^{\prime}\right)  \text{ to }\left(  B^{\prime},B\right)  \right) \\
&  =\left(  \text{\# of nipats from }\left(  A,A^{\prime}\right)  \text{ to
}\left(  B,B^{\prime}\right)  \right) \\
&  \ \ \ \ \ \ \ \ \ \ -\left(  \text{\# of nipats from }\left(  A,A^{\prime
}\right)  \text{ to }\left(  B^{\prime},B\right)  \right)  .
\end{align*}
In view of (\ref{pf.prop.lgv.2paths.count.1}), this rewrites as%
\begin{align*}
&  \det\left(
\begin{array}
[c]{cc}%
\left(  \text{\# of paths from }A\text{ to }B\right)  & \left(  \text{\# of
paths from }A\text{ to }B^{\prime}\right) \\
\left(  \text{\# of paths from }A^{\prime}\text{ to }B\right)  & \left(
\text{\# of paths from }A^{\prime}\text{ to }B^{\prime}\right)
\end{array}
\right) \\
&  =\left(  \text{\# of nipats from }\left(  A,A^{\prime}\right)  \text{ to
}\left(  B,B^{\prime}\right)  \right) \\
&  \ \ \ \ \ \ \ \ \ \ -\left(  \text{\# of nipats from }\left(  A,A^{\prime
}\right)  \text{ to }\left(  B^{\prime},B\right)  \right)  .
\end{align*}
This completes our proof of Proposition \ref{prop.lgv.2paths.count}.
\end{proof}

As in Example \ref{exa.lgv.2paths.2}, the proposition becomes particularly
nice when we have $\left(  \text{\# of nipats from }\left(  A,A^{\prime
}\right)  \text{ to }\left(  B^{\prime},B\right)  \right)  =0$. Here is a
sufficient criterion for when this happens:

\begin{proposition}
[baby Jordan curve theorem]\label{prop.lgv.jordan-2}Let $A$, $B$, $A^{\prime}$
and $B^{\prime}$ be four lattice points satisfying%
\begin{align}
\operatorname*{x}\left(  A^{\prime}\right)   &  \leq\operatorname*{x}\left(
A\right)  ,\ \ \ \ \ \ \ \ \ \ \operatorname*{y}\left(  A^{\prime}\right)
\geq\operatorname*{y}\left(  A\right)  ,\label{eq.prop.lgv.jordan-2.1}\\
\operatorname*{x}\left(  B^{\prime}\right)   &  \leq\operatorname*{x}\left(
B\right)  ,\ \ \ \ \ \ \ \ \ \ \operatorname*{y}\left(  B^{\prime}\right)
\geq\operatorname*{y}\left(  B\right)  . \label{eq.prop.lgv.jordan-2.2}%
\end{align}
Here, $\operatorname*{x}\left(  P\right)  $ and $\operatorname*{y}\left(
P\right)  $ denote the two coordinates of any point $P\in\mathbb{Z}^{2}$.

Let $p$ be any path from $A$ to $B^{\prime}$. Let $p^{\prime}$ be any path
from $A^{\prime}$ to $B$. Then, $p$ and $p^{\prime}$ have a vertex in common.
\end{proposition}

Note that the condition (\ref{eq.prop.lgv.jordan-2.1}) can be restated as
\textquotedblleft the point $A^{\prime}$ lies weakly northwest of
$A$\textquotedblright, where \textquotedblleft weakly northwest of
$A$\textquotedblright\ allows for the options \textquotedblleft due north of
$A$\textquotedblright, \textquotedblleft due west of $A$\textquotedblright%
\ and \textquotedblleft at $A$\textquotedblright. Likewise,
(\ref{eq.prop.lgv.jordan-2.2}) can be restated as \textquotedblleft the point
$B^{\prime}$ lies weakly northwest of $B$\textquotedblright. The following
picture illustrates the situation of Proposition \ref{prop.lgv.jordan-2}:%
\[%
%TCIMACRO{\TeXButton{mse fig 1}{% [tikz figure removed]}}%
%BeginExpansion
% [tikz figure removed]%
%EndExpansion
\]


Proposition \ref{prop.lgv.jordan-2} has an intuitive plausibility to it (one
can think of the path $p$ as creating a \textquotedblleft
river\textquotedblright\ that the path $p^{\prime}$ must necessarily cross
somewhere), but it is not obvious from a mathematical perspective. We give a
rigorous proof in Section \ref{sec.details.det.comb}.

Proposition \ref{prop.lgv.2paths.count} is just the $k=2$ case of a more
general theorem that we will soon derive; however, it already has a nice application:

\begin{corollary}
\label{cor.lgv.binom-unimod}Let $n,k\in\mathbb{N}$. Then, $\dbinom{n}{k}%
^{2}\geq\dbinom{n}{k-1}\cdot\dbinom{n}{k+1}$.
\end{corollary}

\begin{proof}
[Proof of Corollary \ref{cor.lgv.binom-unimod} (sketched).]This is easy to see
algebraically, but here is a combinatorial proof: Define four lattice points
$A=\left(  1,0\right)  $ and $A^{\prime}=\left(  0,1\right)  $ and $B=\left(
k+1,n-k\right)  $ and $B^{\prime}=\left(  k,n-k+1\right)  $. Then, Proposition
\ref{prop.lgv.2paths.count} yields%
\begin{align*}
&  \det\left(
\begin{array}
[c]{cc}%
\left(  \text{\# of paths from }A\text{ to }B\right)  & \left(  \text{\# of
paths from }A\text{ to }B^{\prime}\right) \\
\left(  \text{\# of paths from }A^{\prime}\text{ to }B\right)  & \left(
\text{\# of paths from }A^{\prime}\text{ to }B^{\prime}\right)
\end{array}
\right) \\
&  =\left(  \text{\# of nipats from }\left(  A,A^{\prime}\right)  \text{ to
}\left(  B,B^{\prime}\right)  \right)  -\underbrace{\left(  \text{\# of nipats
from }\left(  A,A^{\prime}\right)  \text{ to }\left(  B^{\prime},B\right)
\right)  }_{\substack{=0\\\text{(since Proposition \ref{prop.lgv.jordan-2}
yields that any path}\\\text{from }A\text{ to }B^{\prime}\text{ and any path
from }A^{\prime}\text{ to }B\text{ have a}\\\text{vertex in common)}}}\\
&  =\left(  \text{\# of nipats from }\left(  A,A^{\prime}\right)  \text{ to
}\left(  B,B^{\prime}\right)  \right)  \geq0.
\end{align*}
However, Proposition \ref{prop.lgv.1-paths.ct} yields
\begin{align*}
&  \left(
\begin{array}
[c]{cc}%
\left(  \text{\# of paths from }A\text{ to }B\right)  & \left(  \text{\# of
paths from }A\text{ to }B^{\prime}\right) \\
\left(  \text{\# of paths from }A^{\prime}\text{ to }B\right)  & \left(
\text{\# of paths from }A^{\prime}\text{ to }B^{\prime}\right)
\end{array}
\right) \\
&  =\left(
\begin{array}
[c]{cc}%
\dbinom{n}{k} & \dbinom{n}{k-1}\\
\dbinom{n}{k+1} & \dbinom{n}{k}%
\end{array}
\right)  ,
\end{align*}
so the determinant on the left hand side is $\dbinom{n}{k}^{2}-\dbinom{n}%
{k-1}\cdot\dbinom{n}{k+1}$. Thus, we have obtained%
\[
\dbinom{n}{k}^{2}-\dbinom{n}{k-1}\cdot\dbinom{n}{k+1}\geq0,
\]
and this proves Corollary \ref{cor.lgv.binom-unimod}.
\end{proof}

Corollary \ref{cor.lgv.binom-unimod} is often stated as \textquotedblleft the
sequence $\dbinom{n}{0},\dbinom{n}{1},\ldots,\dbinom{n}{n}$ is
log-concave\textquotedblright. There are many more log-concave sequences in
combinatorics (see, e.g., \cite{Sagan19}, \cite{Stanle89} and \cite{Brande14}
for more).

\subsubsection{The LGV lemma for $k$ paths}

The following proposition -- which is one of the weakest forms of the
\emph{LGV lemma} (short for \emph{Lindstr\"{o}m--Gessel--Viennot lemma}) --
extends the logic of Proposition \ref{prop.lgv.2paths.count} to nipats between
$k$-vertices for general $k$).

\begin{proposition}
[LGV lemma, lattice counting version]\label{prop.lgv.kpaths.count}Let
$k\in\mathbb{N}$. Let $\mathbf{A}=\left(  A_{1},A_{2},\ldots,A_{k}\right)  $
and $\mathbf{B}=\left(  B_{1},B_{2},\ldots,B_{k}\right)  $ be two
$k$-vertices. Then,%
\begin{align*}
&  \det\left(  \left(  \text{\# of paths from }A_{i}\text{ to }B_{j}\right)
_{1\leq i\leq k,\ 1\leq j\leq k}\right) \\
&  =\sum_{\sigma\in S_{k}}\left(  -1\right)  ^{\sigma}\left(  \text{\# of
nipats from }\mathbf{A}\text{ to }\sigma\left(  \mathbf{B}\right)  \right)  .
\end{align*}

\end{proposition}

The right hand side of this equality can be viewed as a signed count of all
$k$-tuples of paths that start at the points $A_{1},A_{2},\ldots,A_{k}$ in
\textbf{this} order, but end at the points $B_{1},B_{2},\ldots,B_{k}$ in
\textbf{some} order. For example, for $k=3$, the claim of Proposition
\ref{prop.lgv.kpaths.count} takes the form%
\begin{align*}
&  \det\left(  \left(  \text{\# of paths from }A_{i}\text{ to }B_{j}\right)
_{1\leq i\leq3,\ 1\leq j\leq3}\right) \\
&  =\left(  \text{\# of nipats from }\left(  A_{1},A_{2},A_{3}\right)  \text{
to }\left(  B_{1},B_{2},B_{3}\right)  \right) \\
&  \ \ \ \ \ \ \ \ \ \ -\left(  \text{\# of nipats from }\left(  A_{1}%
,A_{2},A_{3}\right)  \text{ to }\left(  B_{1},B_{3},B_{2}\right)  \right) \\
&  \ \ \ \ \ \ \ \ \ \ -\left(  \text{\# of nipats from }\left(  A_{1}%
,A_{2},A_{3}\right)  \text{ to }\left(  B_{2},B_{1},B_{3}\right)  \right) \\
&  \ \ \ \ \ \ \ \ \ \ +\left(  \text{\# of nipats from }\left(  A_{1}%
,A_{2},A_{3}\right)  \text{ to }\left(  B_{2},B_{3},B_{1}\right)  \right) \\
&  \ \ \ \ \ \ \ \ \ \ +\left(  \text{\# of nipats from }\left(  A_{1}%
,A_{2},A_{3}\right)  \text{ to }\left(  B_{3},B_{1},B_{2}\right)  \right) \\
&  \ \ \ \ \ \ \ \ \ \ -\left(  \text{\# of nipats from }\left(  A_{1}%
,A_{2},A_{3}\right)  \text{ to }\left(  B_{3},B_{2},B_{1}\right)  \right)  .
\end{align*}


\begin{proof}
[Proof of Proposition \ref{prop.lgv.kpaths.count} (sketched).]We adapt the
idea of our proof of Proposition \ref{prop.lgv.2paths.count}, but we have to
be more systematic now. Define a set%
\[
\mathcal{A}:=\left\{  \left(  \sigma,\mathbf{p}\right)  \ \mid\ \sigma\in
S_{k}\text{, and}\ \mathbf{p}\text{ is a path tuple from }\mathbf{A}\text{ to
}\sigma\left(  \mathbf{B}\right)  \right\}
\]
\footnote{This set $\mathcal{A}$ is meant to generalize the set $\mathcal{A}$
that was used in the proof of Proposition \ref{prop.lgv.2paths.count}. The
reason why we are defining it to be
\begin{align*}
&  \left\{  \left(  \sigma,\mathbf{p}\right)  \ \mid\ \sigma\in S_{k}\text{,
and}\ \mathbf{p}\text{ is a path tuple from }\mathbf{A}\text{ to }%
\sigma\left(  \mathbf{B}\right)  \right\} \\
&  \text{instead of }\left\{  \mathbf{p}\ \mid\ \mathbf{p}\text{ is a path
tuple from }\mathbf{A}\text{ to }\sigma\left(  \mathbf{B}\right)  \text{ for
some }\sigma\in S_{k}\right\}
\end{align*}
is to make sure that each path tuple in $\mathcal{A}$ \textquotedblleft
remembers\textquotedblright\ which permutation $\sigma\in S_{k}$ it comes
from. (This is the same rationale that caused us to take the disjoint union in
the proof of Proposition \ref{prop.lgv.2paths.count}; but now we are
explicitly inserting the $\sigma$ into the elements of $\mathcal{A}$ rather
than handwaving about disjoint unions.)}. Define a subset $\mathcal{X}$ of
$\mathcal{A}$ by\footnote{Here we are saying that a pair $\left(
\sigma,\mathbf{p}\right)  \in\mathcal{A}$ is an \emph{ipat} if the path tuple
$\mathbf{p}$ is an ipat, and we are saying that a pair $\left(  \sigma
,\mathbf{p}\right)  \in\mathcal{A}$ is a \emph{nipat} if the path tuple
$\mathbf{p}$ is a nipat. This is a bit sloppy ($\sigma$ has nothing to do with
whether $\mathbf{p}$ is an ipat or a nipat), but we hope that no confusion
will ensue.}%
\[
\mathcal{X}:=\left\{  \text{ipats in }\mathcal{A}\right\}  =\left\{  \left(
\sigma,\mathbf{p}\right)  \in\mathcal{A}\ \mid\ \mathbf{p}\text{ is
intersecting}\right\}  .
\]
Set%
\[
\operatorname*{sign}\left(  \sigma,\mathbf{p}\right)  :=\left(  -1\right)
^{\sigma}\ \ \ \ \ \ \ \ \ \ \text{for each }\left(  \sigma,\mathbf{p}\right)
\in\mathcal{A}.
\]


Again, we need to find a sign-reversing involution $f:\mathcal{X}%
\rightarrow\mathcal{X}$ that has no fixed points.

Again, we construct this involution by exchanging the tails of two
intersecting paths\footnote{This is probably obvious, but just in case: We say
that two paths \emph{intersect} if they have a vertex in common.} in our path
tuple. There is a complication now, due to the fact that there might be
several pairs of intersecting paths. We have to come up with a rule for
picking one such pair so that when we apply $f$ again to the result of the
exchange, then we again pick the same pair. Otherwise, $f$ won't be an involution!

There are different ways to do this. Here is one: If $\left(  \sigma,\left(
p_{1},p_{2},\ldots,p_{k}\right)  \right)  \in\mathcal{X}$, then we construct
$f\left(  \sigma,\left(  p_{1},p_{2},\ldots,p_{k}\right)  \right)
\in\mathcal{X}$ as follows:

\begin{itemize}
\item We say that a point $u$ is \emph{crowded} if it is contained in at least
two of our paths $p_{1},p_{2},\ldots,p_{k}$. Since $\left(  p_{1},p_{2}%
,\ldots,p_{k}\right)  $ is intersecting, there exists at least one crowded point.

\item We pick the smallest $i\in\left[  k\right]  $ such that $p_{i}$ contains
a crowded point.

\item Then, we pick the first crowded point $v$ on $p_{i}$.

\item Then, we pick the largest $j\in\left[  k\right]  $ such that $v$ belongs
to $p_{j}$. (Note that $j>i$, since $v$ is crowded.)

\item Call the part of $p_{i}$ that comes after $v$ the \emph{tail} of $p_{i}%
$. Call the part of $p_{i}$ that comes before $v$ the \emph{head}\textbf{ }of
$p_{i}$.

Call the part of $p_{j}$ that comes after $v$ the \emph{tail} of $p_{j}$. Call
the part of $p_{j}$ that comes before $v$ the \emph{head}\textbf{ }of $p_{j}$.

\item Then, we exchange the tails of the paths $p_{i}$ and $p_{j}$ (while
leaving all other paths unchanged).

\item We let $\left(  q_{1},q_{2},\ldots,q_{k}\right)  $ be the resulting path
tuple\footnote{Thus,%
\[
q_{i}=\left(  \text{head of }p_{i}\right)  \cup\left(  \text{tail of }%
p_{j}\right)  \ \ \ \ \ \ \ \ \ \ \text{and}\ \ \ \ \ \ \ \ \ \ q_{j}=\left(
\text{head of }p_{j}\right)  \cup\left(  \text{tail of }p_{i}\right)
\]
(where \textquotedblleft head\textquotedblright\ means \textquotedblleft part
until $v$\textquotedblright, and \textquotedblleft tail\textquotedblright%
\ means \textquotedblleft part after $v$\textquotedblright). Furthermore, we
have $q_{m}=p_{m}$ for any $m\in\left[  k\right]  \setminus\left\{
i,j\right\}  $, since we have left all paths other than $p_{i}$ and $p_{j}$
unchanged.}.

\item We set $\sigma^{\prime}:=\sigma t_{i,j}$, where $t_{i,j}$ is the
transposition in $S_{k}$ that swaps $i$ with $j$. Thus, $\left(  q_{1}%
,q_{2},\ldots,q_{k}\right)  $ is a path tuple from $\mathbf{A}$ to
$\sigma^{\prime}\left(  \mathbf{B}\right)  $ (because exchanging the tails of
the paths $p_{i}$ and $p_{j}$ has switched their ending points $B_{\sigma
\left(  i\right)  }$ and $B_{\sigma\left(  j\right)  }$ to $B_{\sigma\left(
j\right)  }=B_{\sigma^{\prime}\left(  i\right)  }$ and $B_{\sigma\left(
i\right)  }=B_{\sigma^{\prime}\left(  j\right)  }$, respectively).

\item Finally, we set%
\[
f\left(  \sigma,\left(  p_{1},p_{2},\ldots,p_{k}\right)  \right)  :=\left(
\sigma^{\prime},\left(  q_{1},q_{2},\ldots,q_{k}\right)  \right)  .
\]

\end{itemize}

This defines a map $f:\mathcal{X}\rightarrow\mathcal{X}$ (again, it is not
hard to see that it is well-defined). Here is an example: If $k=5$ and if
$\left(  p_{1},p_{2},\ldots,p_{k}\right)  $ is%
\[%
%TCIMACRO{\TeXButton{five paths}{% [tikz figure removed]}}%
%BeginExpansion
% [tikz figure removed]%
%EndExpansion
\]
(where $B_{m}^{\prime}$ is shorthand for $B_{\sigma\left(  m\right)  }$), then
$v$ is the point where paths $p_{2}$ and $p_{4}$ intersect, and we have $i=2$
and $j=4$, and therefore $\left(  q_{1},q_{2},\ldots,q_{k}\right)  $ is%
\[%
%TCIMACRO{\TeXButton{five paths}{% [tikz figure removed]}}%
%BeginExpansion
% [tikz figure removed]%
%EndExpansion
\]
(where we again have drawn the exchanged tails extra-thick).

Convince yourself that the map $f$ defined above really is a sign-reversing
involution from $\mathcal{X}$ to $\mathcal{X}$. (This means showing that
applying $f$ twice in succession to any given $\left(  \sigma,\mathbf{p}%
\right)  \in\mathcal{X}$ returns $\left(  \sigma,\mathbf{p}\right)  $, and
that $\operatorname*{sign}\left(  f\left(  \sigma,\mathbf{p}\right)  \right)
=-\operatorname*{sign}\left(  \sigma,\mathbf{p}\right)  $ for any $\left(
\sigma,\mathbf{p}\right)  \in\mathcal{X}$. The proof of the second claim, of
course, relies on parts \textbf{(b)} and \textbf{(d)} of Proposition
\ref{prop.perm.sign.props}.)

Thus, we have defined a sign-reversing involution $f:\mathcal{X}%
\rightarrow\mathcal{X}$. This involution $f$ has no fixed points (since it is
sign-reversing). It is now easy to complete the proof: Lemma
\ref{lem.sign.cancel2} yields%
\[
\sum_{\left(  \sigma,\mathbf{p}\right)  \in\mathcal{A}}\operatorname*{sign}%
\left(  \sigma,\mathbf{p}\right)  =\sum_{\left(  \sigma,\mathbf{p}\right)
\in\mathcal{A}\setminus\mathcal{X}}\operatorname*{sign}\left(  \sigma
,\mathbf{p}\right)  .
\]
In view of our definition of $\operatorname*{sign}\left(  \sigma
,\mathbf{p}\right)  $, this rewrites as%
\begin{equation}
\sum_{\left(  \sigma,\mathbf{p}\right)  \in\mathcal{A}}\left(  -1\right)
^{\sigma}=\sum_{\left(  \sigma,\mathbf{p}\right)  \in\mathcal{A}%
\setminus\mathcal{X}}\left(  -1\right)  ^{\sigma}.
\label{pf.prop.lgv.kpaths.count.at}%
\end{equation}
The left hand side of this equality is%
\begin{align*}
\sum_{\left(  \sigma,\mathbf{p}\right)  \in\mathcal{A}}\left(  -1\right)
^{\sigma}  &  =\sum_{\sigma\in S_{k}}\left(  -1\right)  ^{\sigma
}\underbrace{\left(  \text{\# of path tuples from }\mathbf{A}\text{ to }%
\sigma\left(  \mathbf{B}\right)  \right)  }_{\substack{=\prod_{i=1}^{k}\left(
\text{\# of paths from }A_{i}\text{ to }B_{\sigma\left(  i\right)  }\right)
\\\text{(by the product rule, since a path tuple from }\mathbf{A}\text{ to
}\sigma\left(  \mathbf{B}\right)  \text{ is just}\\\text{a tuple }\left(
p_{1},p_{2},\ldots,p_{k}\right)  \text{, where each }p_{i}\text{ is a path
from }A_{i}\text{ to }B_{\sigma\left(  i\right)  }\text{)}}}\\
&  \ \ \ \ \ \ \ \ \ \ \ \ \ \ \ \ \ \ \ \ \left(  \text{by the definition of
}\mathcal{A}\right) \\
&  =\sum_{\sigma\in S_{k}}\left(  -1\right)  ^{\sigma}\prod_{i=1}^{k}\left(
\text{\# of paths from }A_{i}\text{ to }B_{\sigma\left(  i\right)  }\right) \\
&  =\det\left(  \left(  \text{\# of paths from }A_{i}\text{ to }B_{j}\right)
_{1\leq i\leq k,\ 1\leq j\leq k}\right)
\end{align*}
(by the definition of the determinant), whereas the right hand side is%
\begin{align*}
\sum_{\left(  \sigma,\mathbf{p}\right)  \in\mathcal{A}\setminus\mathcal{X}%
}\left(  -1\right)  ^{\sigma}  &  =\sum_{\left(  \sigma,\mathbf{p}\right)
\in\mathcal{A}\text{ is a nipat}}\left(  -1\right)  ^{\sigma}%
\ \ \ \ \ \ \ \ \ \ \left(
\begin{array}
[c]{c}%
\text{since }\mathcal{X}=\left\{  \text{ipats in }\mathcal{A}\right\} \\
\text{entails }\mathcal{A}\setminus\mathcal{X}=\left\{  \text{nipats in
}\mathcal{A}\right\}
\end{array}
\right) \\
&  =\sum_{\sigma\in S_{k}}\left(  -1\right)  ^{\sigma}\left(  \text{\# of
nipats from }\mathbf{A}\text{ to }\sigma\left(  \mathbf{B}\right)  \right)
\end{align*}
(by the definition of $\mathcal{A}$). Thus, (\ref{pf.prop.lgv.kpaths.count.at}%
) rewrites as
\begin{align*}
&  \det\left(  \left(  \text{\# of paths from }A_{i}\text{ to }B_{j}\right)
_{1\leq i\leq k,\ 1\leq j\leq k}\right) \\
&  =\sum_{\sigma\in S_{k}}\left(  -1\right)  ^{\sigma}\left(  \text{\# of
nipats from }\mathbf{A}\text{ to }\sigma\left(  \mathbf{B}\right)  \right)  .
\end{align*}
This proves Proposition \ref{prop.lgv.kpaths.count}.
\end{proof}

